\documentclass[12pt]{article}

\usepackage[utf8]{inputenc}
\usepackage[T1]{fontenc}
\usepackage[spanish]{babel}
\usepackage{amsmath}
\usepackage{amssymb}
\usepackage{xcolor}
\usepackage{geometry}

\geometry{
    a4paper,
    left=1.5cm,
    right=1.5cm,
    top=1.3cm,
    bottom=1.5cm
}

\setlength{\parindent}{0pt}
\pagestyle{empty}

\begin{document}

% ============================================================
% PAUTA PROBLEMA 1
% ============================================================

\section*{Pauta Problema 1}

Considerando

\[
V=k\int\frac{dq}{r}
\]

\subsection*{a) \hfill \textcolor{red}{\textbf{(12 pts)}}}

El potencial en el origen está dado por

\[
V_O=V_A+V_B
\qquad
\textcolor{red}{\textbf{1 pt}}
\]

Para cada caso es necesario encontrar $dq$, $r$ y los límites de integración.

\vspace{0.3cm}

\textbf{Arco de circunferencia}

\[
dq_A=\lambda_0R\,d\theta
\qquad
\textcolor{red}{\textbf{1 pt}}
\]

\[
r_{AO}=R
\qquad
\textcolor{red}{\textbf{1 pt}}
\]

\[
\frac{3\pi}{2}
\longrightarrow
2\pi+\frac{\pi}{4}
=
\frac{9\pi}{4}
\qquad
\textcolor{red}{\textbf{1 pt}}
\]

\vspace{0.3cm}

\textbf{Barra}

\[
dq_B
=
\lambda(x)\,dx
=
\lambda_1\frac{x}{R}\,dx
\qquad
\textcolor{red}{\textbf{1 pt}}
\]

\[
r_{BO}=x
\qquad
\textcolor{red}{\textbf{1 pt}}
\]

\[
R\longrightarrow3R
\qquad
\textcolor{red}{\textbf{1 pt}}
\]

\vspace{0.4cm}

\textbf{Reemplazando}

Para el arco,

\[
V_A
=
\int_{\frac{3\pi}{2}}^{\frac{9\pi}{4}}
k\frac{dq_A}{r_{AO}}
=
\int_{\frac{3\pi}{2}}^{\frac{9\pi}{4}}
k\frac{\lambda_0R}{R}\,d\theta
\qquad
\textcolor{red}{\textbf{1 pt}}
\]

\[
V_A
=
k\lambda_0
\int_{\frac{3\pi}{2}}^{\frac{9\pi}{4}}
d\theta
=
k\lambda_0
\left(
\frac{9\pi}{4}-\frac{3\pi}{2}
\right)
\]

\[
\boxed{
V_A
=
k\lambda_0\frac{3\pi}{4}
}
\qquad
\textcolor{red}{\textbf{1 pt}}
\]

Para la barra,

\[
V_B
=
\int_R^{3R}
k\frac{\lambda_1\dfrac{x}{R}}{x}\,dx
=
\frac{k\lambda_1}{R}
\int_R^{3R}dx
\qquad
\textcolor{red}{\textbf{1 pt}}
\]

\[
V_B
=
\frac{k\lambda_1}{R}(3R-R)
\]

\[
\boxed{
V_B=2k\lambda_1
}
\qquad
\textcolor{red}{\textbf{1 pt}}
\]

Por lo tanto,

\[
\boxed{
V_O
=
k\lambda_0\frac{3\pi}{4}
+
2k\lambda_1
}
\qquad
\textcolor{red}{\textbf{1 pt}}
\]

% ============================================================

\subsection*{b) \hfill \textcolor{red}{\textbf{(3 pts)}}}

Considerando la condición $V_O=0$,

\[
V_O
=
k\left(
\frac{3\pi}{4}\lambda_0+2\lambda_1
\right)
=
0
\qquad
\textcolor{red}{\textbf{1 pt}}
\]

\[
\frac{3\pi}{4}\lambda_0+2\lambda_1=0
\qquad
\textcolor{red}{\textbf{1 pt}}
\]

Por lo tanto,

\[
\boxed{
\lambda_0
=
-\frac{8}{3\pi}\lambda_1
}
\qquad
\textcolor{red}{\textbf{1 pt}}
\]

% ============================================================

\subsection*{c) \hfill \textcolor{red}{\textbf{(5 pts)}}}

La carga total está dada por

\[
q_{\text{tot}}
=
q_A+q_B
\qquad
\textcolor{red}{\textbf{1 pt}}
\]

Para el arco,

\[
q_A
=
\lambda_0R
\int_{\frac{3\pi}{2}}^{\frac{9\pi}{4}}d\theta
=
\lambda_0R\frac{3\pi}{4}
\qquad
\textcolor{red}{\textbf{1 pt}}
\]

Utilizando

\[
\lambda_0
=
-\frac{8}{3\pi}\lambda_1,
\]

se obtiene

\[
q_A
=
\left(
-\frac{8}{3\pi}\lambda_1
\right)
R\frac{3\pi}{4}
\]

\[
\boxed{
q_A=-2\lambda_1R
}
\qquad
\textcolor{red}{\textbf{1 pt}}
\]

Para la barra,

\[
q_B
=
\frac{\lambda_1}{R}
\int_R^{3R}x\,dx
=
\frac{\lambda_1}{R}
\left[
\frac{x^2}{2}
\right]_R^{3R}
\]

\[
q_B
=
\frac{\lambda_1}{R}
\left(
\frac{9R^2-R^2}{2}
\right)
\]

\[
\boxed{
q_B=4\lambda_1R
}
\qquad
\textcolor{red}{\textbf{1 pt}}
\]

Finalmente,

\[
q_{\text{tot}}
=
-2\lambda_1R+4\lambda_1R
\]

\[
\boxed{
q_{\text{tot}}
=
2\lambda_1R
}
\qquad
\textcolor{red}{\textbf{1 pt}}
\]

\newpage

% ============================================================
% PAUTA PROBLEMA 2
% ============================================================

\section*{Pauta Problema 2}

Considerando

\[
V=k\int\frac{dq}{r}
\]

\subsection*{a) \hfill \textcolor{red}{\textbf{(12 pts)}}}

El potencial en el origen está dado por

\[
V_O=V_A+V_B
\qquad
\textcolor{red}{\textbf{1 pt}}
\]

Para cada caso es necesario encontrar $dq$, $r$ y los límites de integración.

\vspace{0.3cm}

\textbf{Arco de circunferencia}

\[
dq_A
=
\lambda(\theta)R\,d\theta
=
\lambda_1\cos\theta\,R\,d\theta
\qquad
\textcolor{red}{\textbf{1 pt}}
\]

\[
r_{AO}=R
\qquad
\textcolor{red}{\textbf{1 pt}}
\]

\[
\frac{3\pi}{4}
\longrightarrow
\frac{5\pi}{4}
\qquad
\textcolor{red}{\textbf{1 pt}}
\]

\vspace{0.3cm}

\textbf{Barra}

Aunque la barra se encuentra sobre el eje $x$ negativo, $x$ representa
la distancia entre el origen y cada elemento de carga, por lo que se
considera positiva.

\[
dq_B=\lambda_0\,dx
\qquad
\textcolor{red}{\textbf{1 pt}}
\]

\[
r_{BO}=x
\qquad
\textcolor{red}{\textbf{1 pt}}
\]

\[
R\longrightarrow3R
\qquad
\textcolor{red}{\textbf{1 pt}}
\]

\vspace{0.4cm}

\textbf{Reemplazando}

Para el arco,

\[
V_A
=
\int_{\frac{3\pi}{4}}^{\frac{5\pi}{4}}
k\frac{\lambda_1\cos\theta\,R}{R}\,d\theta
\qquad
\textcolor{red}{\textbf{1 pt}}
\]

\[
V_A
=
k\lambda_1
\left[
\sin\theta
\right]_{\frac{3\pi}{4}}^{\frac{5\pi}{4}}
\]

\[
V_A
=
k\lambda_1
\left(
-\frac{\sqrt{2}}{2}
-
\frac{\sqrt{2}}{2}
\right)
\]

\[
\boxed{
V_A=-\sqrt{2}\,k\lambda_1
}
\qquad
\textcolor{red}{\textbf{1 pt}}
\]

Para la barra,

\[
V_B
=
\int_R^{3R}
k\frac{\lambda_0}{x}\,dx
=
k\lambda_0
\int_R^{3R}\frac{dx}{x}
\qquad
\textcolor{red}{\textbf{1 pt}}
\]

\[
V_B
=
k\lambda_0
\left[
\ln x
\right]_R^{3R}
=
k\lambda_0
\ln\left(\frac{3R}{R}\right)
\]

\[
\boxed{
V_B=k\lambda_0\ln3
}
\qquad
\textcolor{red}{\textbf{1 pt}}
\]

Por lo tanto,

\[
\boxed{
V_O
=
k\left(
\lambda_0\ln3-\sqrt{2}\lambda_1
\right)
}
\qquad
\textcolor{red}{\textbf{1 pt}}
\]

% ============================================================

\subsection*{b) \hfill \textcolor{red}{\textbf{(3 pts)}}}

Considerando la condición $V_O=0$,

\[
V_O
=
k\left(
\lambda_0\ln3-\sqrt{2}\lambda_1
\right)
=
0
\qquad
\textcolor{red}{\textbf{1 pt}}
\]

\[
\lambda_0\ln3-\sqrt{2}\lambda_1=0
\qquad
\textcolor{red}{\textbf{1 pt}}
\]

Por lo tanto,

\[
\boxed{
\lambda_0
=
\frac{\sqrt{2}}{\ln3}\lambda_1
}
\qquad
\textcolor{red}{\textbf{1 pt}}
\]

% ============================================================

\subsection*{c) \hfill \textcolor{red}{\textbf{(5 pts)}}}

La carga total está dada por

\[
q_{\text{tot}}
=
q_A+q_B
\qquad
\textcolor{red}{\textbf{1 pt}}
\]

Para el arco,

\[
q_A
=
\lambda_1R
\int_{\frac{3\pi}{4}}^{\frac{5\pi}{4}}
\cos\theta\,d\theta
\]

\[
q_A
=
\lambda_1R
\left[
\sin\theta
\right]_{\frac{3\pi}{4}}^{\frac{5\pi}{4}}
\]

\[
\boxed{
q_A=-\sqrt{2}\lambda_1R
}
\qquad
\textcolor{red}{\textbf{1 pt}}
\]

Para la barra,

\[
q_B
=
\lambda_0
\int_R^{3R}dx
=
\lambda_0(3R-R)
\]

\[
\boxed{
q_B=2\lambda_0R
}
\qquad
\textcolor{red}{\textbf{1 pt}}
\]

Utilizando

\[
\lambda_0
=
\frac{\sqrt{2}}{\ln3}\lambda_1,
\]

se obtiene

\[
\boxed{
q_B
=
\frac{2\sqrt{2}}{\ln3}\lambda_1R
}
\qquad
\textcolor{red}{\textbf{1 pt}}
\]

Finalmente,

\[
q_{\text{tot}}
=
-\sqrt{2}\lambda_1R
+
\frac{2\sqrt{2}}{\ln3}\lambda_1R
\]

\[
\boxed{
q_{\text{tot}}
=
\sqrt{2}\lambda_1R
\left(
\frac{2}{\ln3}-1
\right)
}
\qquad
\textcolor{red}{\textbf{1 pt}}
\]

\end{document}